% Documento generado desde la guía DOCX equivalente.
% Compilación recomendada: pdflatex Guia_XX_*.tex (dos pasadas).
\documentclass[11pt,a4paper]{article}
\usepackage[utf8]{inputenc}
\usepackage[T1]{fontenc}
\usepackage[spanish,es-nodecimaldot]{babel}
\usepackage[a4paper,margin=2.35cm]{geometry}
\usepackage{lmodern}
\usepackage{microtype}
\usepackage[table]{xcolor}
\usepackage{array}
\usepackage{tabularx}
\usepackage{float}
\usepackage{enumitem}
\usepackage[most]{tcolorbox}
\usepackage{fancyhdr}
\usepackage{hyperref}
\definecolor{uemcgreen}{HTML}{004C3F}
\definecolor{uemcgold}{HTML}{E5B93F}
\definecolor{tfgblue}{HTML}{20566B}
\definecolor{tablelight}{HTML}{E7F1EE}
\definecolor{uemclight}{HTML}{F2F7F5}
\hypersetup{colorlinks=true,linkcolor=uemcgreen,urlcolor=tfgblue,pdftitle={Tecnologías y decisiones},pdfauthor={Alejandro Villarrubia García}}
\setlength{\parindent}{0pt}
\setlength{\parskip}{6pt}
\renewcommand{\arraystretch}{1.16}
\setlist[itemize]{leftmargin=1.2cm,itemsep=2pt,topsep=2pt}
\setlist[enumerate]{leftmargin=1.2cm,itemsep=4pt,topsep=2pt}
\pagestyle{fancy}
\fancyhf{}
\fancyhead[L]{\small\color{uemcgreen}GUÍA 02 · DECISIONES TÉCNICAS}
\fancyhead[R]{\small\color{uemcgreen}TFG · UEMC}
\fancyfoot[C]{\color{uemcgreen}\thepage}
\renewcommand{\headrulewidth}{0.4pt}
\renewcommand{\headrule}{\hbox to\headwidth{\color{uemcgold}\leaders\hrule height \headrulewidth\hfill}}
\begin{document}
\begin{titlepage}
  \centering
  \vspace*{2.4cm}
  {\Large\color{uemcgreen}\textbf{GUÍA 02 · DECISIONES TÉCNICAS}\par}
  \vspace{1.5cm}
  {\Huge\bfseries\color{uemcgreen}Tecnologías y decisiones\par}
  \vspace{0.7cm}
  {\large\color{tfgblue}Cómo funciona cada elección y cómo justificarla frente a alternativas\par}
  \vfill
  {\large Proyecto TFG · versión alineada con el código actual\par}
  \vspace{0.35cm}
  {\large Alejandro Villarrubia García\par}
  \vspace{0.35cm}
  {\large Universidad Europea Miguel de Cervantes\par}
\end{titlepage}
\tableofcontents
\clearpage

\section{Criterio de selección}

La arquitectura prioriza reproducibilidad local, facilidad de defensa, separación de responsabilidades y coste operativo bajo. Las decisiones no intentan convertir el TFG en un servicio de producción multiusuario: delimitan un detector demostrable, auditable y extensible.

\section{Comparativa de tecnologías}

\begin{table}[H]
  \centering
  \small
  \caption{Comparativa de tecnologías}
  \label{tab:comparativa-1}
  \begin{tabularx}{\textwidth}{|>{\raggedright\arraybackslash}p{2.45cm}|X|X|X|}
    \hline
    \rowcolor{uemcgreen}
    \textcolor{white}{\textbf{Tecnología}} & \textcolor{white}{\textbf{Cómo se usa}} & \textcolor{white}{\textbf{Por qué se elige}} & \textcolor{white}{\textbf{Alternativa y límite}} \\
    \hline
    Python & Integra parser, reglas, ML, API y automatizaciones & Ecosistema científico, lectura clara y buena trazabilidad & Java/Node serían válidos; añadirían complejidad para este alcance \\
    \hline
    \rowcolor{tablelight}
    Streamlit & Cliente web con sistema visual compartido y adaptable & Permite validar el flujo cliente-servidor con una interfaz coherente y poco código & React daría comunicación navegador-API directa, pero exigiría otro stack \\
    \hline
    scikit-learn & TF-IDF vectoriza y MLPClassifier clasifica & Pipeline reproducible, API madura y métricas estándar & TensorFlow sería más flexible para deep learning, pero sobredimensionado para texto tabular \\
    \hline
    \rowcolor{tablelight}
    TF-IDF & Convierte palabras y bigramas en variables & Explicable, rápido y adecuado para corpus moderados & Transformers mejorarían contexto, con más coste y menor explicabilidad local \\
    \hline
    email estándar & Parsea MIME y cabeceras EML & No añade dependencia para la parte esencial del formato & Librerías externas pueden aportar atajos, pero no son necesarias \\
    \hline
    \rowcolor{tablelight}
    BeautifulSoup & Inspecciona HTML y formularios & Tolera HTML imperfecto de correos & Parser regex puro sería frágil; un navegador sería inseguro e innecesario \\
    \hline
    Gmail API + OAuth & Obtiene mensajes con consentimiento del usuario & Permisos oficiales y flujo revocable & IMAP sería más genérico, pero no aporta el mismo control API \\
    \hline
    \rowcolor{tablelight}
    HTTP stdlib & Expone análisis y administración central & Cero framework extra, contrato pequeño y fácil de auditar & FastAPI facilitaría OpenAPI y escalado; no es necesario para el prototipo \\
    \hline
    joblib & Persiste pipeline TF-IDF + MLP & Guarda vectorizador y modelo como una unidad & ONNX sería más portable; exige conversión y no elimina el requisito de confiar en artefactos \\
    \hline
    \rowcolor{tablelight}
    unittest + Ruff & Prueba comportamiento y calidad estática & Incluidos en Python y rápidos en CI/local & pytest aporta fixtures más ricas; el conjunto actual es suficiente y directo \\
    \hline
  \end{tabularx}
\end{table}

\section{Por qué dos modelos lingüísticos}

El detector identifica idioma por mensaje, no por sesión. Mantiene un modelo español y otro inglés porque las palabras, stopwords y corpus tienen distribuciones distintas. EmailAnalysisService cachea un detector por idioma para combinar corrección con rendimiento.

\begin{tcolorbox}[colback=uemclight,colframe=uemcgreen,title=Nota para la defensa]
\textbf{Respuesta breve:} La alternativa de un único modelo multilingüe simplificaría ficheros, pero mezclaría vocabularios y haría más difícil justificar qué datos explican cada predicción.
\end{tcolorbox}

\section{SOLID aplicado}

\begin{table}[H]
  \centering
  \small
  \caption{Aplicación de SOLID}
  \label{tab:comparativa-3}
  \begin{tabularx}{\textwidth}{|>{\raggedright\arraybackslash}p{2.75cm}|X|X|}
    \hline
    \rowcolor{uemcgreen}
    \textcolor{white}{\textbf{Principio}} & \textcolor{white}{\textbf{Aplicación en el código}} & \textcolor{white}{\textbf{Beneficio demostrable}} \\
    \hline
    Responsabilidad única & Parser, señales, scorer, explicaciones, ML, monitor y transporte están separados & Un cambio en URLs no obliga a tocar Gmail o Streamlit \\
    \hline
    \rowcolor{tablelight}
    Abierto/cerrado & EmailAnalysisService recibe analyzer, detector loader y detector de idioma inyectables & Se prueban estrategias y se añaden modelos sin reescribir el coordinador \\
    \hline
    Sustitución & Los consumidores dependen de BackendClient o su protocolo mínimo & Web, extensión y monitor consumen el mismo contrato HTTP \\
    \hline
    \rowcolor{tablelight}
    Segregación & El almacenamiento, el análisis y el transporte tienen superficies pequeñas & Cada prueba necesita solo la interfaz que utiliza \\
    \hline
    Inversión & NeuralModelTrainer depende de ModelStorage, no de una ruta fija & Facilita memoria, disco y dobles de prueba \\
    \hline
  \end{tabularx}
\end{table}

\section{Decisiones de seguridad}

\begin{itemize}
  \item El parser limita EML a 10 MiB; /analyze admite 16 MiB y las rutas de datasets 256 MiB con Content-Length obligatorio.
  \item La API no acepta CORS *; permite Gmail y orígenes chrome-extension://, y no devuelve trazas internas.
  \item Por defecto Streamlit y el backend escuchan en loopback. Para una demo móvil puede abrirse solo Streamlit en 0.0.0.0:8501 y mantener el backend en 127.0.0.1:8766; no requiere --allow-remote, no es Internet y debe limitarse a una LAN privada de confianza. Un backend remoto sí exige HTTPS y controles adicionales.
  \item El modelo no ejecuta archivos adjuntos; solo registra metadatos y señales.
  \item Los nuevos joblib no conservan textos brutos de entrenamiento. Los modelos antiguos se sanean al cargarse y guardarse con el formato actual.
  \item Credenciales, tokens, estados y artefactos temporales de QA se mantienen fuera del índice Git mediante .gitignore.
\end{itemize}

\section{Reproducibilidad y evaluación}

Los hiperparámetros finales son TF-IDF (ngram\_range 1-2, max\_features 3.000, min\_df 1, strip\_accents unicode) y MLP ((64, 32), relu, alpha 0,0001, learning rate 0,001, 500 iteraciones, sin early stopping, semilla 42). El protocolo usa 1.148/209 textos ES y 65.661/16.416 EN. La clase positiva mezcla spam/phishing según la fuente y se documenta como proxy, no como equivalencia. Los hashes permiten identificar los joblib y los modelos no guardan textos brutos.

El benchmark reproducible separa arranque e inferencia. Se midió en Windows 11 Home 25H2 de 64 bits (26200.9278), Ryzen 7 7800X3D, 63,1 GiB RAM y Python 3.12.7 con scikit-learn 1.9.0, Streamlit 1.58.0, NumPy 2.4.6, SciPy 1.17.1 y joblib 1.5.3. La carga diferida redujo importación heurística un 96,6 \% y arranque de aplicación un 76,2 \%; la inferencia varió menos de ±3 \%.

La calibración usa 40 casos separados en cinco particiones: pesos heurísticos 20-50, umbrales 20-60 y alta confianza 65-85. Selecciona 45/55, umbral 21 y alta confianza 70 con desempates predefinidos; no produce una probabilidad calibrada. La prueba de concurrencia usa el backend real: 8 y 16 análisis terminaron sin errores, pero a 32 aparecieron fallos. Se recomiendan 4-8 para demo.

La validación automatizada reúne 94 pruebas Python y 2 pruebas con Chromium. Incluye idioma determinista y el recorrido local Gmail EML -> JSON -> HTTP -> backend -> Telegram. GitHub Actions instala versiones fijadas, ejecuta Ruff, verifica calibración, regenera la evaluación, comprueba el respaldo de defensa y recorre web-backend. OAuth real usa una lista E2E separada porque sus credenciales no pueden entrar en CI.

\section{Diferencias con la propuesta inicial}

\begin{table}[H]
  \centering
  \small
  \caption{Diferencias frente a la propuesta inicial}
  \label{tab:comparativa-4}
  \begin{tabularx}{\textwidth}{|>{\raggedright\arraybackslash}p{2.75cm}|X|X|}
    \hline
    \rowcolor{uemcgreen}
    \textcolor{white}{\textbf{Propuesta}} & \textcolor{white}{\textbf{Implementación actual}} & \textcolor{white}{\textbf{Cómo explicarlo}} \\
    \hline
    TensorFlow & scikit-learn MLP & Se eligió una red suficiente para el corpus y más reproducible/ligera en el entorno académico \\
    \hline
    \rowcolor{tablelight}
    Blacklists/reputación & Señales locales de dominio, URL, punycode y acortadores & Se evita depender de red externa y se mantiene privacidad; reputación online queda como extensión \\
    \hline
    Certificados & No se inspecciona TLS desde el correo & El EML no prueba por sí mismo el certificado del destino; se prioriza análisis estático seguro \\
    \hline
    \rowcolor{tablelight}
    Interfaz simple & Streamlit como cliente de una API central obligatoria & La interfaz sigue siendo simple y todos los canales comparten los modelos del servidor \\
    \hline
  \end{tabularx}
\end{table}

\section{Qué no prometer}

\begin{itemize}
  \item No decir que se consulta una blacklist online si no se ha configurado una fuente externa.
  \item No presentar el accuracy de entrenamiento como precisión en producción.
  \item No llamar a la extensión un producto publicado: se carga en modo desarrollador y consume el backend central configurado.
\end{itemize}

\end{document}
